\documentclass[spanish,12pt,a4paper]{book}% Preambulo

%Todos los comandos de Latex se escriben en \ ya que es el simbolo de comandos
%Dentro del comando \documentclass[]{} posemos agregar 
%[Aqui podemos agregar condiciones adicionales de la libreria]
%{El tipo de documento que queremos generar}
\usepackage{hyperref}
\usepackage{geometry}
\usepackage{amssymb}
\usepackage{babel}
\geometry{letterpaper,margin=1.5cm}


\begin {document}% Cuerpo del documento
\title{Primeros pasos}

{\Large Hablemos de la funcion ''usepackege''}\\\\
\begin{enumerate}
\item \textbf{{Paquetes}}\\\\
	 Los paquetes que veremos seran:
	\begin{itemize}
		\item modern e inputenc:  {\small (Estos se usan para mejorar el manejo de caracteres especiales o para que se pueda copiar y pegar texto correctamente del pdf generado)}
		\item babel:\  {\small (Este es un diccionario que advierte al documento que idioma usaras al escribir)}
		\item mathtools:\  {\small (Ayuda a la escritura matematica)}
		\item hyperref:\  {\small (Nos ayuda a incluir ligas para en el documento)} 
		\item biblatex:\  {\small (Para administrar la bibliografia)}
		\item tiks:\  {\small (Para agregar todo tipo de ilustraciones)}
	\end{itemize}

\item \textbf{{Portada}}\\\\
	Los siguientes son comandos que afectaran al incio de la portada:\label{Tomamos el capitulo o en este caso el item}\\\\
	\textbackslash title\\ \textbackslash author\\ \textbackslash date 

\item\textbf{{Texto}}\\\\
	Los saltos de linea se pueden hacer con el comando:\\\\
	\textbackslash \textbackslash \\
	\textbackslash newline\\
	\\
	Los saltos de pagina se dan por:\\\\
	\textbackslash newpage\\
	\textbackslash pagebreak (Aqui solicitamos un salto de pagina)\\
	\\
	Cuando nos interesa meter distintos tipos de letras ocuparemos :\\
	\\
	\textbackslash textit (cursivas)
	\\
	\textbackslash textbf (negritas)
	\\
	\textbackslash textsc (versalilla)
	\\	
	\textsc{Esta es la versalilla}
	\\
	\\
	\newpage
	Los comandos mas comunes para cambiar el tamaño de letra son:
	\\\\
	\textbackslash tiny\\\textbackslash scriptsize\\\textbackslash small\\ \textbackslash normalisize\\ \textbackslash huge\\ \textbackslash Huge
\item \textbf{Listas :}\\
	Los comandos de listas se hacen generando un inicio begin y su final end, existen dos tipos.
	\begin{itemize}
		\item \textbackslash itemize
	\end{itemize}
	\begin{enumerate}
		\item \textbackslash enumerate (se generan letras ya que usamos para todo el texto un enumerate)
	\end{enumerate}
\item \textbf{Imagenes}
	Si deseamos integrar imagenes usaremos carpetas incluidas de tal forma que el archivo tenga
	\begin{enumerate}
		\item Carpeta con las imagenes a incluir
		\item El documento del texto
		\item El auxiliar del documento
	\end{enumerate}
	Es necesario que en el preambulo agreguemos el userpage 'graphicx', de esta forma podremos agregar imagenes usando el comando 'includegraphicx'\\
	Existe el entorno 'figure' y el comando titulo 'caption' con el que podremos agregar un titulo a la figura por ejemplo:
	\begin{itemize}
		\item \textbackslash begin \{figure\} $\left[ h \right]$ 
		\item \textbackslash centering
		\item \textbackslash includegraphics [width = 0.4 \textbackslash textwidth]\{Nombre de la carpeta / Nombre de la foto .jpg\} 
		\item \textbackslash caption \{Texto de la figura\}
		\item \textbackslash label \{Ver marcador del pdf\}
		\item \textbackslash end \{figure\}
	\end{itemize}
	En los corchetes de figure te indica donde quieres poner la foto:\\
	\textbf{h}(here), \textbf{t}(arriba), \textbf{b}(abajo) o \textbf{p}(para juntar las imagenes en una pagina)
	\\ Para ajustar altura usamos \textbf{'height'} y ancho \textbf{'width'}\\
	Para escalar usaremos \textbf{'scale' } donde usaremos la escala y sus grados\\\\
	Para imagenes a un lado del texto usaremos el paquete \textbf{''wrapfig''}, esto bajo el entorno\\
	 \textbf{''wrapfigure''}:
	\begin{itemize}
		\item \textbackslash begin \{wrapfigure\}\{Lado\}\{Porcentage del tamaño\textbf{\textbackslash textwidth}\}
		\item \textbackslash includegraphics \{{Ancho \textbackslash textwidth}\}\{Foto\}
		\item \textbackslash end \{wrapfigure\}
	\end{itemize}
\newpage
\item \textbf{Hipertextos}\\\\
	Es normal querer hacer mencion o citar cosa dentro del mismo documento para eso usaremos los hipertextos y etiquetas:
	\\\\
	Usaremos la paqueteria\textbf{ \textbackslash usepackage 'hyperreg'}y con el comando \textbf{\textbackslash label 'nombre-enlace'} definiremos el nombre del enlace dentro del documento al cual se podra hacer referencia posteriormente mediante\textbf{ \textbackslash ref}
	\\\\Añadire al documento dicha libreria y referenciare cualquier cosa para verla aqui:\ref{Tomamos el capitulo o en este caso el item}	\\\\
	Tambien se puede hacer referencia de una figura o tabla, si se desea agregar URL ent usaremos los comandos \textbf{\textbackslash url} o si se desea poner un texto alternativo se usara el comando \textbf{\textbackslash href}
	\\\\
\item \textbf{Texto matematico} 
	\begin{itemize}
		\item \$ \$ si solo deseamo agregar texto matematico
		\item \$\$ \$\$ si se desea agregar texto centrado matematico
	\end{itemize}
	Para generar ecuaciones numeradad utilizaremos el entorno ' equation ' de la siguiengte forma:
	\begin{itemize}	
	\item \textbackslash begin equation
	\item La ecuacion que deseamos poner
	\item \textbackslash end equation
	\end{itemize}	
	Para generar letras de conjuntos del campo usaremos la paqueteria \textbf{\textbackslash usepackage\{amssymb\}}\\
	Para generar funciones usaremos:\\\\
	 T:\textbf{\textbackslash mathbb\{R\}2\textbackslash longmapsto\textbackslash mathbb\{R\}2} = $(T:\mathbb{R}^2\longmapsto\mathbb{R}^2)$
\item \textbf{Tablas}\\\\
	Generalmente las tablas se hacen mediante los entornos ' table ' y 'tabular':
	\begin{itemize}
		\item \textbackslash begin \{table\}
		\item centering (para centrar el texto de la tabla)
		\item \textbackslash begin \{tabular\} \{Izquierda Centro Derecha\} Se puede agregar $ |c|$ para rectas verticales
		\item Si queremos ponerle una linea en toda la tabla de manera horizontal usamos \textbf{\textbackslash hline}
		\item \textbackslash end \{tabular\}
		\item \textbackslash end \{table\}
	\end{itemize}
	Usalmente las tablas no siempre generan el mejor tamñao por lo que usaremos algunos paquetes para ello\\
	El paquete graphicx incluye el comando \textbf{\textbackslash 'resizebox'} , el cual nos sirve en el entorno \textbf{'table'} y asigna un tamaño especifico
	\\
\newpage
\item \textbf{Listas de figuras y tablas}\\\\
	Se puede obtener fácilmente las listas de figuras y tablas en el documento con los comandos \textbf{\textbackslash listoffigures} y \textbf{\textbackslash listoftables} respectivamente\\
\item \textbf{Bibliografia}\\\\
 	Existen dos tipo de entornos para citar en Latex:\\
	\begin{enumerate}
 		\item thebibilography\\
 		Este sirve para generar referencias utilizadas  y citarlas en el texto. Este entorno generalmente se localiza al final del texto, en la seccióm hecha especialmente para la bibliografia y este es su esqueleto:\\
 		\begin{itemize}
 			\item \textbackslash begin \{thebibliography\}{Aqui el numero max de citas a poner}
 			\item \textbackslash begin \{El nombre que le pondremos a la cita\} Seguido la bibliografia \textbf{(No acepta el formato APA)}
 			\item (Podemos poner una cantidad finita a la cantidad de citas que pusimos en el entorno)
 			\item  \textbackslash end \{thebibliography\}
 		\end{itemize}
 		Para citar en el texto usaremos el comando \textbf{\textbackslash cite \{Aqui el nombre unico que se le puso a la cita\}}
 		\item Bibtex
    	\\Para la forma de citar con este entorno revisar el marcador del Pdf
 	\end{enumerate}
\item \textbf{Tamaño del documento}\\\\
	Para cambiar el tamaño del documento ocuparemos la paqueteria \textbf{\textbackslash geometry} y el comando \textbf{\textbackslash geometry \{Formato del documento, Los margenes (''left'', ''right'', ''top'', ''bottom'') o ''magin'' seguido de los sentimetros\}}, esto en el preambulo\\
	Lo pondremos en el formato \textbf{letterpaper, margin = 2cm}\\
\item \textbf{Indice}\\\\
	Para generar un indice usaremos el comando \textbf{\textbackslash tableofcontens}, esto va en el cuerpo del documento


\end{enumerate}
\end {document}
